\section{Alternative Views}
\label{sec:alternatives}

\paragraph{Memory is a sixth primitive.} Memory is prominent in agent architectures \citep{sumers2024coala,packer2023memgpt}, and could be argued to deserve its own place. We place it instead as a store in the world that \prim{Context} reads and \prim{Output} writes, governed by the persistence concern. The model has no memory of its own, and every implementation we examined (memory files loaded into the instructions, stores searched through tools, summaries carried forward) is a decision about what \prim{Context} includes and which actions \prim{Output} offers. Making memory a primitive would split one decision, what the model is shown, across two components.

\paragraph{Planning is a sixth primitive.} Planning is either something the model does in its own tokens, which needs no harness support, or a structure the harness imposes: a planning mode that restricts actions, a plan kept in a file, a separate planning call. Each of these is \prim{Control flow} (which steps run, in what order, under what restrictions) or \prim{Context} (where the plan is kept and shown).

\paragraph{Multi-agent systems are something else.} A system of several agents, each a model in a harness, is itself a harness whose \prim{Output} includes delegating to another agent and whose \prim{Input} includes that agent's result. Sub-agents, hand-offs between agents and orchestrator--worker patterns are \prim{Control flow} decisions at this outer level. The primitives compose: a harness of harnesses still has five.

\paragraph{The harness will be absorbed into the model.} Models increasingly learn behaviours that harnesses once imposed, such as deciding when to stop, calling tools in a trained format, and reasoning before acting. But a model cannot open a file, keep a log, or refuse a command on its own: whatever is learned, something must still carry tokens to the world and back. Training moves \emph{policy} into the model (what to do); the \emph{mechanism} (how anything happens) remains in the harness. If anything, the primitives become more important as models are trained inside a particular harness (\S\ref{sec:implications}), because the harness then defines the world the model was trained for.

\paragraph{The primitives are too coarse to be useful.} Five boxes can describe anything, and so might explain nothing. Our answer is that the primitives are a coordinate system, not an explanation: their value is in locating a design decision, a failure or a difference between two systems in one place, so that it can be studied, compared and changed independently of the rest. The design spaces in \S\ref{sec:primitives} and the case study in \S\ref{sec:casestudy} are where the content lies.
